\section{Introducción}
\textit{Responsable: Mariano (consolida el informe).}

Nubeduca es un LMS (\textit{Learning Management System}) ofrecido como SaaS multi-tenant: una sola
plataforma que usan a la vez varias instituciones educativas (\texttt{utec}, \texttt{upc} y
\texttt{pucp}), cada una con sus datos aislados. Los docentes crean cursos, matriculan alumnos,
publican módulos y tareas, califican y publican notas; los estudiantes descargan materiales,
entregan tareas, ven sus notas y reciben avisos por correo.

La solución es \textbf{serverless} (API Gateway, Lambda y DynamoDB) y \textbf{orientada a eventos}
(EventBridge, SQS y SNS). Se divide en seis microservicios, cada uno con su API, sus funciones Lambda
y su propia tabla DynamoDB, y se despliega íntegramente con Serverless Framework en AWS Academy
Learner Lab.

\subsection{Objetivos}
\begin{itemize}
  \item Cumplir los requisitos del enunciado: multi-tenancy, serverless, arquitectura basada en
        eventos, mínimo cinco microservicios protegidos, base de datos propia por microservicio con su
        estructura JSON documentada, 30\,000 registros por microservicio en tres tenants, uso de S3, SNS
        y SQS, despliegue con Serverless Framework y frontend en Amplify.
  \item Demostrar de punta a punta, con dos tenants, el flujo académico del curso piloto:
        contenido $\rightarrow$ tarea $\rightarrow$ entrega $\rightarrow$ nota $\rightarrow$ aviso.
\end{itemize}

\subsection{Alcance}
El proyecto entrega la Fase 0 (fundaciones de plataforma) y la Fase 1 (MVP académico). Las fases 2 a
4 (quizzes, foros, analítica, integraciones externas) quedan como \textit{roadmap}.

\subsection{Cómo leer este informe}
La sección~\ref{sec:arquitectura} presenta la arquitectura; la~\ref{sec:microservicios}, cada
microservicio con la estructura JSON de su tabla; la~\ref{sec:eventos}, el bus de eventos; la
\ref{sec:seguridad}, la autenticación y el aislamiento por tenant; la~\ref{sec:frontend}, el frontend;
la~\ref{sec:despliegue}, el despliegue y la observabilidad; la~\ref{sec:pruebas}, las pruebas; la
\ref{sec:limitaciones}, las limitaciones del Learner Lab, y la~\ref{sec:repos}, los repositorios. El
anexo relaciona cada requisito del enunciado con la sección que lo cubre.
